\subsection{The Three Outcomes and a Simple Parameter Example}

Every linear system has exactly one solution, no solution, or infinitely many solutions. For a square system $Ax=b$, a nonzero determinant guarantees the first case. A zero determinant removes uniqueness but does not decide existence.

\begin{remarkbox}
The condition $\det A=0$ is a warning sign, not a complete classification.
\end{remarkbox}

Consider
\[
\begin{cases}
x+y=1,\\
2x+2y=a.
\end{cases}
\]
Twice the first equation gives $2x+2y=2$.

\begin{summarybox}[Classification]
If $a=2$, the equations describe the same line and there are infinitely many solutions. If $a\neq2$, they describe parallel distinct lines and there is no solution.
\end{summarybox}
